A prop de la Lluna hi ha un camp elèctric que, en la cara il·luminada, està dirigit cap a l'exterior de la Lluna i, en la cara fosca, cap al centre. Tot i que a la Lluna no hi ha atmosfera, aquests camps elèctrics poden mantenir partícules de pols en suspensió.

\figura[0.5]{fig1.png}

A la superfície de la cara il·luminada, el mòdul del camp és $10\ \mathrm{N\,C^{-1}}$, mentre que a la superfície de la cara fosca és $1,0\ \mathrm{N\,C^{-1}}$.

\begin{apartat}{1}
Calculeu la relació $\dfrac{q}{m}$ (càrrega elèctrica/massa) que ha de tenir una partícula de pols situada a la cara il·luminada de la Lluna perquè es trobi en situació d'equilibri de forces. Expliciteu el signe que ha de tenir la càrrega elèctrica.
\end{apartat}

\begin{apartat}{1}
Considereu una partícula amb una càrrega $q = +10\ \mathrm{nC}$ i una massa $m = 0,020\ \mathrm{mg}$ situada a la cara fosca de la Lluna. Calculeu la força total que actua sobre la partícula i el temps que tardarà a recórrer 10 metres partint del repòs.
\end{apartat}

\blocetiquetat{Dada:}{$g\,(\text{Lluna}) = 1,62\ \mathrm{m\,s^{-2}}$.}
